% GENERATED FILE. Edit canonical lessons, then run npm run generate:books.
% canonical-source-hash: fnv1a64:65429e5f07c54043
% canonical-lessons: FR-C55-joie, FR-C55-peur, FR-C55-colere, FR-C55-tristesse, FR-C55-heureux

\chapter{Joy, Fear, Anger, Sadness, Happy}
\label{ch:fr-joy-fear-anger-sadness-happy}
\hlchaptermodality{55}

\begin{chapteropening}
\textbf{By the end of this chapter:} \emph{I can name joy, fear, anger, sadness, happy.}

Complete the last lesson of chapter 55: I can name joy, fear, anger, sadness, happy.
\end{chapteropening}

\section[la joie]{la joie — joy}

\label{lesson:FR-C55-joie}



\begin{quote}
Before the new one: say the French for tiredness, then the French for ill.
\end{quote}

\subsection*{You'll want to know: la joie}
\textbf{la joie} — \textquotedblleft{}joy\textquotedblright{}.

\begin{cousinweb}
\textbf{joie} $\leftarrow$ Latin \emph{gaudia}, \textquotedblleft{}joys\textquotedblright{}, the plural of \emph{gaudium} (EtymDB 2.1) — inherited: the \emph{g} became \emph{j}, as in \emph{jambe}.
\end{cousinweb}

\begin{grammarlens}[title={What you've built}]
One more word for this chapter.
\end{grammarlens}

\subsection*{Your turn}
\begin{itemize}
\raggedright
  \item \emph{Say it:} \emph{la joie}
  \item \emph{Say it:} \emph{la joie}, once more
  \item \emph{Say it:} say \emph{la joie}
  \item \emph{Recall:} say the French for tiredness, then the French for ill, then say \emph{la joie} again
\end{itemize}

\begin{tcolorbox}[breakable,colback=teal!4,colframe=teal!35!black,title={Before you move on}]
What does la joie mean? (\textquotedblleft{}Joy\textquotedblright{}.) Say it once more.
\end{tcolorbox}
\section[la peur]{la peur — fear}

\label{lesson:FR-C55-peur}



\begin{quote}
Before the new one: say the French for ill, then the French for joy.
\end{quote}

\subsection*{You'll want to know: la peur}
\textbf{la peur} — \textquotedblleft{}fear\textquotedblright{}. \textbf{j'ai peur}: I am afraid.

\begin{cousinweb}
\textbf{peur} $\leftarrow$ Old French \emph{peor} $\leftarrow$ Latin \emph{pavor}, \textquotedblleft{}fear, fright\textquotedblright{} (EtymDB 2.1) — inherited; the \emph{v} between vowels dropped out.
\end{cousinweb}

\begin{grammarlens}[title={What you've built}]
One more word for this chapter.
\end{grammarlens}

\subsection*{Your turn}
\begin{itemize}
\raggedright
  \item \emph{Say it:} \emph{la peur}
  \item \emph{Say it:} \emph{la peur}, once more
  \item \emph{Say it:} say \emph{la peur}
  \item \emph{Recall:} say the French for ill, then the French for joy, then say \emph{la peur} again
\end{itemize}

\begin{tcolorbox}[breakable,colback=teal!4,colframe=teal!35!black,title={Before you move on}]
What does la peur mean? (\textquotedblleft{}Fear\textquotedblright{}.) Say it once more.
\end{tcolorbox}
\section[la colère]{la colère — anger}

\label{lesson:FR-C55-colere}



\begin{quote}
Before the new one: say the French for joy, then the French for fear.
\end{quote}

\subsection*{You'll want to know: la colère}
\textbf{la colère} — \textquotedblleft{}anger\textquotedblright{}.

\begin{cousinweb}
\textbf{colère} $\leftarrow$ Latin \emph{cholera}, \textquotedblleft{}a bilious illness\textquotedblright{}, from Greek \emph{kholera}, built on \emph{khole}, \textquotedblleft{}bile\textquotedblright{} — a learned borrowing: anger was blamed on bile. (EtymDB 2.1 traces Italian \emph{collera} to the same Latin word.)
\end{cousinweb}

\begin{grammarlens}[title={What you've built}]
One more word for this chapter.
\end{grammarlens}

\subsection*{Your turn}
\begin{itemize}
\raggedright
  \item \emph{Say it:} \emph{la colère}
  \item \emph{Say it:} \emph{la colère}, once more
  \item \emph{Say it:} say \emph{la colère}
  \item \emph{Recall:} say the French for joy, then the French for fear, then say \emph{la colère} again
\end{itemize}

\begin{tcolorbox}[breakable,colback=teal!4,colframe=teal!35!black,title={Before you move on}]
What does la colère mean? (\textquotedblleft{}Anger\textquotedblright{}.) Say it once more.
\end{tcolorbox}
\section[la tristesse]{la tristesse — sadness}

\label{lesson:FR-C55-tristesse}



\begin{quote}
Before the new one: say the French for fear, then the French for anger.
\end{quote}

\subsection*{You'll want to know: la tristesse}
\textbf{la tristesse} — \textquotedblleft{}sadness\textquotedblright{}.

\begin{cousinweb}
\textbf{tristesse} $\leftarrow$ Latin \emph{tristitia}, \textquotedblleft{}sadness\textquotedblright{}, from \emph{tristis}, \textquotedblleft{}sad\textquotedblright{} — inherited; Latin \emph{-itia} became French \emph{-esse}. (EtymDB 2.1 records the Latin word.)
\end{cousinweb}

\begin{grammarlens}[title={What you've built}]
One more word for this chapter.
\end{grammarlens}

\subsection*{Your turn}
\begin{itemize}
\raggedright
  \item \emph{Say it:} \emph{la tristesse}
  \item \emph{Say it:} \emph{la tristesse}, once more
  \item \emph{Say it:} say \emph{la tristesse}
  \item \emph{Recall:} say the French for fear, then the French for anger, then say \emph{la tristesse} again
\end{itemize}

\begin{tcolorbox}[breakable,colback=teal!4,colframe=teal!35!black,title={Before you move on}]
What does la tristesse mean? (\textquotedblleft{}Sadness\textquotedblright{}.) Say it once more.
\end{tcolorbox}
\section[heureux]{heureux — happy}

\label{lesson:FR-C55-heureux}



\begin{quote}
Before the new one: say the French for anger, then the French for sadness.
\end{quote}

\subsection*{You'll want to know: heureux}
\textbf{heureux} — \textquotedblleft{}happy\textquotedblright{}.

\begin{cousinweb}
\textbf{heureux} — built from Old French \emph{eur}, \textquotedblleft{}luck\textquotedblright{}, with \emph{-eux}, \textquotedblleft{}full of\textquotedblright{}: lucky, then happy. \emph{eur} $\leftarrow$ Latin \emph{augurium}, \textquotedblleft{}an omen\textquotedblright{}. (EtymDB 2.1)
\end{cousinweb}

\begin{grammarlens}[title={What you've built}]
One more word for this chapter.
\end{grammarlens}

\subsection*{Your turn}
\begin{itemize}
\raggedright
  \item \emph{Say it:} \emph{heureux}
  \item \emph{Say it:} \emph{heureux}, once more
  \item \emph{Say it:} say \emph{heureux}
  \item \emph{Recall:} say the French for anger, then the French for sadness, then say \emph{heureux} again
\end{itemize}

\begin{tcolorbox}[breakable,colback=teal!4,colframe=teal!35!black,title={Before you move on}]
What does heureux mean? (\textquotedblleft{}Happy\textquotedblright{}.) Say it once more.
\end{tcolorbox}
